\textbf{What are the principles of Cloud Native architecture?}

Cloud-native architecture represents a paradigm shift in software development and deployment, emphasizing the creation of applications that are designed to leverage cloud computing's inherent advantages. The principles of cloud-native architecture can be categorized into several key areas: microservices, containerization, scalability, resilience, and continuous delivery.

At the core of cloud-native architecture is the microservices approach, which decomposes applications into smaller, loosely coupled services that can be developed, deployed, and scaled independently. This architectural style enhances flexibility and allows for faster development cycles, as teams can work on different services simultaneously without impacting the entire application (Oyeniran, 2024; Kazanavicius & Mažeika, 2023; Balalaie et al., 2016). Microservices also facilitate the use of diverse technology stacks, enabling organizations to select the best tools for each service, which can lead to improved performance and maintainability (Souza et al., 2019).

Containerization is another fundamental principle of cloud-native architecture. By encapsulating microservices within containers, developers can ensure that applications run consistently across various environments, from development to production (Hasan et al., 2023; Shi, 2023). Containers provide an isolated environment for each service, allowing for easier management and deployment while maximizing resource utilization (Toffetti et al., 2017). This approach also supports the implementation of Infrastructure as Code (IaC), which automates the provisioning and management of infrastructure, further enhancing operational efficiency (Oyeniran, 2024; Shi, 2023).

Scalability is a critical aspect of cloud-native architecture, enabling applications to handle varying loads by dynamically adjusting resources. This is achieved through auto-scaling mechanisms that automatically increase or decrease the number of active instances of a service based on current demand (Shi, 2023; Henning & Hasselbring, 2022). The microservices architecture inherently supports scalability, as individual services can be scaled independently, allowing organizations to optimize resource allocation and reduce costs (Souza et al., 2019; Goshi et al., 2021).

Resilience is another essential principle, as cloud-native applications are designed to withstand failures and recover quickly. This is often accomplished through redundancy, where multiple instances of a service are deployed across different locations, ensuring that if one instance fails, others can take over (Oyeniran, 2024; Duan, 2021). Additionally, cloud-native applications often incorporate self-healing capabilities, allowing them to automatically detect and recover from failures without human intervention (Toffetti et al., 2017).

Finally, continuous delivery is a vital principle in cloud-native architecture, promoting frequent and reliable software releases. This is facilitated by practices such as continuous integration and continuous deployment (CI/CD), which automate the testing and deployment processes, enabling teams to deliver new features and updates rapidly (Oyeniran, 2024; Shi, 2023; Sebrechts et al., 2022). By adopting CI/CD practices, organizations can respond quickly to market changes and user feedback, maintaining a competitive edge in their respective industries (Kazanavicius & Mažeika, 2023).

In summary, the principles of cloud-native architecture—microservices, containerization, scalability, resilience, and continuous delivery—collectively enable organizations to build and operate applications that are flexible, efficient, and capable of adapting to changing demands. These principles not only enhance the development process but also improve the overall performance and reliability of applications in a cloud environment.

References:
Balalaie, A., Heydarnoori, A., & Jamshidi, P. (2016). Migrating to cloud-native architectures using microservices: an experience report., 201-215. https://doi.org/10.1007/978-3-319-33313-7_15
Duan, Q. (2021). Intelligent and autonomous management in cloud-native future networks—a survey on related standards from an architectural perspective. Future Internet, 13(2), 42. https://doi.org/10.3390/fi13020042
Goshi, E., Jarschel, M., Pries, R., He, M., & Kellerer, W. (2021). Investigating inter-nf dependencies in cloud-native 5g core networks.. https://doi.org/10.23919/cnsm52442.2021.9615565
Hasan, M., Osman, M., Admodisastro, N., & Muhammad, S. (2023). A quality driven framework for decomposing legacy monolith applications to microservice architecture.. https://doi.org/10.21203/rs.3.rs-3060410/v1
Henning, S. and Hasselbring, W. (2022). A configurable method for benchmarking scalability of cloud-native applications. Empirical Software Engineering, 27(6). https://doi.org/10.1007/s10664-022-10162-1
Kazanavicius, J. and Mažeika, D. (2023). Evaluation of microservice communication while decomposing monoliths. Computing and Informatics, 42(1), 1-36. https://doi.org/10.31577/cai_2023_1_1
Oyeniran, O. (2024). A comprehensive review of leveraging cloud-native technologies for scalability and resilience in software development. International Journal of Science and Research Archive, 11(2), 330-337. https://doi.org/10.30574/ijsra.2024.11.2.0432
Sebrechts, M., Volckaert, B., Turck, F., Yang, K., & Al-Naday, M. (2022). Fog native architecture: intent-based workflows to take cloud native toward the edge. Ieee Communications Magazine, 60(8), 44-50. https://doi.org/10.1109/mcom.003.2101075
Shi, H. (2023). Comedge: cloud-native platform for integrated computing and communication in satellite–terrestrial network. Electronics, 12(20), 4252. https://doi.org/10.3390/electronics12204252
Souza, P., Marques, W., Oliveira, R., & Ferreto, T. (2019). Iagree: infrastructure-agnostic resilience benchmark tool for cloud native platforms.. https://doi.org/10.5220/0007728503960403
Toffetti, G., Brunner, S., Blöchlinger, M., Spillner, J., & Bohnert, T. (2017). Self-managing cloud-native applications: design, implementation, and experience. Future Generation Computer Systems, 72, 165-179. https://doi.org/10.1016/j.future.2016.09.002